\EMhold

\EMsection{The Laurent Series}{The Laurent Series}

\begin{EMtheorem}{Laurent series}

If the function \EMim{f(z)} is analytic in the domain \EMim{D}, its Laurent series about the point \EMim{z_0} is written
\EMdisp{f(z)\EMbr =\sum_{n=-\infty}^{\infty}a_n\,(z-z_0)^n,\qquad\EMbr  a_n\EMbr =\frac1{2\pi i}\oint_C\frac{f(v)}{(v-z_0)^{n+1}}\,dv}
where \EMim{C} is a circle centred at \EMim{z_0} and the function is analytic at all points inside and on this circle, except possibly at \EMim{z_0}.

The radius of convergence is the distance from \EMim{z_0} to the nearest non-analytic point of \EMim{f(z)}.

\end{EMtheorem}

For non-negative \EMim{n} we have
\EMdisp{a_n\EMbr =\frac1{2\pi i}\oint_C\frac{f(v)}{(v-z_0)^{n+1}}\,dv\EMbr =\frac{f^{(n)}(z_0)}{n!},\qquad\EMbr  n\EMbr =0,1,2,\dots}
Therefore
\EMdisp{\begin{aligned}
f(z)&=\sum_{n=0}^{\infty}\frac{f^{(n)}(z_0)}{n!}(z-z_0)^n+\sum_{n=-\infty}^{-1}a_n(z-z_0)^n\\
&=\sum_{n=0}^{\infty}\frac{f^{(n)}(z_0)}{n!}(z-z_0)^n+\sum_{n=1}^{\infty}a_{-n}(z-z_0)^{-n}\\
&=\sum_{n=0}^{\infty}\frac{f^{(n)}(z_0)}{n!}(z-z_0)^n+\sum_{n=1}^{\infty}\frac{A_n}{(z-z_0)^n},\qquad A_n=a_{-n}
\end{aligned}}
\EMdisp{A_n\EMbr =a_{-n}\EMbr =\frac1{2\pi i}\oint_C\frac{f(v)}{(v-z_0)^{-n+1}}\,dv\EMbr =\frac1{2\pi i}\oint_C(v-z_0)^{n-1}f(v)\,dv,\qquad\EMbr  n\EMbr =1,2,\dots}

\begin{EMimportant}{The Laurent series: Taylor part and principal part}

\EMdisp{f(z)\EMbr =\underbrace{\sum_{n=0}^{\infty}\frac{f^{(n)}(z_0)}{n!}(z-z_0)^n}_{\text{Taylor series part}}+\underbrace{\sum_{n=1}^{\infty}\frac{A_n}{(z-z_0)^n}}_{\text{principal part}},\qquad\EMbr  A_n\EMbr =\frac1{2\pi i}\oint_C(v-z_0)^{n-1}f(v)\,dv}

\end{EMimportant}

If \EMim{z_0} is an irregular (singular) point, in other words if the function is not analytic there, the Taylor part alone is not sufficient to describe the function.

\EMsubsection{The Different Cases for the Point \EMim{z_0}}{The Different Cases for the Point z0}

\begin{enumerate}
\def\labelenumi{\arabic{enumi}.}
\tightlist
\item
  If all the \EMim{A_n} are zero, the point \EMim{z_0} is called a \textbf{regular} point. In this case the Laurent series is just the Taylor series.
\item
  If \EMim{A_p\neq0} but \EMim{(A_n=0\ \ \forall n>p)}, then \EMim{z_0} is called a \textbf{pole of order \EMim{p}}.
\item
  If \EMim{z_0} is not a regular point and is also not a pole (in other words, no \EMim{p} can be found such that \EMim{A_p\neq0} and \EMim{(A_n=0\ \ \forall n>p)}), then \EMim{z_0} is called an \textbf{essential singularity}.
\end{enumerate}

\EMhold

\EMsection{The Residue}{The Residue}

\begin{EMdefinition}{Residue}

The residue of \EMim{f(z)} at the point \EMim{z_0} is \EMim{A_1}.
\EMdisp{\operatorname{Res}\{f(z)\}\Big|_{z=z_0}\EMbr =A_1\EMbr =\frac1{2\pi i}\oint_Cf(v)\,dv}

\end{EMdefinition}

\begin{EMtheorem}{}

If the function \EMim{f(z)} has a pole of order \EMim{p} at \EMim{z_0}, then
\EMdisp{A_k\EMbr =\frac1{(p-k)!}\left.\frac{d^{\,p-k}}{dz^{\,p-k}}\Big[(z-z_0)^p\,f(z)\Big]\right|_{z=z_0},\qquad\EMbr  k\EMbr =1,2,\dots,p}

\end{EMtheorem}

\begin{EMproof}

\EMdisp{f(z)\EMbr =\sum_{n=1}^{\infty}\frac{A_n}{(z-z_0)^n}+\sum_{n=0}^{\infty}\frac{f^{(n)}(z_0)}{n!}(z-z_0)^n}
If \EMim{f(z)} has a pole of order \EMim{p} at \EMim{z_0} we have
\EMdisp{f(z)\EMbr =\frac{A_p}{(z-z_0)^p}+\frac{A_{p-1}}{(z-z_0)^{p-1}}+\cdots+\frac{A_2}{(z-z_0)^2}+\frac{A_1}{z-z_0}+\sum_{n=0}^{\infty}\frac{f^{(n)}(z_0)}{n!}(z-z_0)^n}
Multiplying both sides by \EMim{(z-z_0)^p} gives
\EMdisp{(z-z_0)^pf(z)\EMbr =A_p+A_{p-1}(z-z_0)+\cdots+A_2(z-z_0)^{p-2}+A_1(z-z_0)^{p-1}+\sum_{n=0}^{\infty}\frac{f^{(n)}(z_0)}{n!}(z-z_0)^{n+p}}
Now if we set \EMim{z=z_0} we get
\EMdisp{(z-z_0)^pf(z)\Big|_{z=z_0}\EMbr =A_p}
If we differentiate \EMim{(z-z_0)^pf(z)} once and then set \EMim{z=z_0}:
\EMdisp{\begin{aligned}
\frac{d}{dz}\Big[(z-z_0)^pf(z)\Big]\Big|_{z=z_0}&=\Bigg\{A_{p-1}+\cdots+A_2(p-2)(z-z_0)^{p-3}+A_1(p-1)(z-z_0)^{p-2}\\
&\qquad+\sum_{n=0}^{\infty}\frac{f^{(n)}(z_0)}{n!}(z-z_0)^{n+p-1}(n+p)\Bigg\}\Bigg|_{z=z_0}=A_{p-1}
\end{aligned}}
If we differentiate \EMim{(z-z_0)^pf(z)} twice and then set \EMim{z=z_0}:
\EMdisp{\begin{aligned}
\frac{d^2}{dz^2}\Big[(z-z_0)^pf(z)\Big]\Big|_{z=z_0}&=\Bigg\{2A_{p-2}+\cdots+A_2(p-2)(p-3)(z-z_0)^{p-4}+A_1(p-1)(p-2)(z-z_0)^{p-3}\\
&\qquad+\sum_{n=0}^{\infty}\frac{f^{(n)}(z_0)}{n!}(z-z_0)^{n+p-2}(n+p)(n+p-1)\Bigg\}\Bigg|_{z=z_0}=2A_{p-2}
\end{aligned}}
\EMdisp{\Longrightarrow\quad\EMbr  A_{p-2}\EMbr =\frac1{2!}\left.\frac{d^2}{dz^2}\Big[(z-z_0)^pf(z)\Big]\right|_{z=z_0}}

\end{EMproof}

\begin{EMexercise}{}

Continue the differentiation in the same way and complete the proof of the theorem.

\end{EMexercise}

\begin{EMremark}{Consequence}

In particular, when \EMim{f(z)} has a pole of order \EMim{p} at \EMim{z_0}, the residue is given by
\EMdisp{\operatorname{Res}\{f(z)\}\EMbr =A_1\EMbr =\frac1{(p-1)!}\left.\frac{d^{\,p-1}}{dz^{\,p-1}}\Big[(z-z_0)^p\,f(z)\Big]\right|_{z=z_0}}

\end{EMremark}

\EMhold

\EMsection{The Residue Theorem}{The Residue Theorem}

\begin{EMtheorem}{The residue theorem}

If the function \EMim{f(z)} is analytic everywhere inside the path \EMim{C} except at the points \EMim{z_1,z_2,\dots,z_n}, and the residues of the function at these points are \EMim{A_1^1,A_1^2,\dots,A_1^n} respectively, then
\EMdisp{\oint_Cf(z)\,dz\EMbr =2\pi i\sum_{k=1}^{n}A_1^k}

\end{EMtheorem}

\EMfigure{m06-residue-theorem}{A path \EMim{C} and the singular points \EMim{z_1,\dots,z_n} inside it;
a small circle is drawn around each point}

The proof is left to the reader as an exercise.

\EMhold

\EMsection{Examples of Taylor and Laurent Expansions}{Examples of Taylor and Laurent Expansions}

\begin{EMexample}{}

Find the Laurent series of \EMim{f(z)=\dfrac1{z^4+1}} in a neighbourhood of \EMim{z_0=0}.

We see that \EMim{f(z)} is analytic at \EMim{z_0=0}, so this is a regular point and the Laurent series is the Taylor (Maclaurin) series. Hence
\EMdisp{f(z)\EMbr =f(0)+\frac1{1!}f^{(1)}(0)z+\frac1{2!}f^{(2)}(0)z^2+\frac1{3!}f^{(3)}(0)z^3+\cdots}
Geometric series:
\EMdisp{1+q+q^2+q^3+q^4+\cdots\EMbr =\frac1{1-q},\qquad\EMbr  |q|<1}
\EMdisp{\frac1{1+z^4}\EMbr =1-z^4+z^8-z^{12}+-\cdots,\qquad\EMbr  |z^4|<1\ \EMbr \Longrightarrow\ |z|<1}
(Taylor/Maclaurin series; region of convergence of radius 1.)
\EMdisp{z^4+1\EMbr =0\;\EMbr \Longrightarrow\;z_1\EMbr =e^{i\pi/4},\quad\EMbr  z_2\EMbr =e^{i3\pi/4},\quad\EMbr  z_3\EMbr =e^{i5\pi/4},\quad\EMbr  z_4\EMbr =e^{i7\pi/4}}
We see that the distance from \EMim{z_0=0} to the nearest non-analytic point of the function (the radius of convergence) is \EMim{1}.

\end{EMexample}

\EMfigure{m06-radius-1}{The four poles of \EMim{z^4+1=0} on the unit circle; the radius of
convergence of the series about the origin is \EMim{1}}

\begin{EMexample}{}

Find the Laurent series of \EMim{f(z)=\dfrac{4z^2+30z+68}{(z+4)^2(z-2)}} in a neighbourhood of \EMim{z_0=0}.

We see that \EMim{f(z)} is analytic at \EMim{z_0=0}, so this is a regular point and the Laurent series is the Taylor (Maclaurin) series. Hence
\EMdisp{f(z)\EMbr =f(0)+\frac1{1!}f^{(1)}(0)z+\frac1{2!}f^{(2)}(0)z^2+\frac1{3!}f^{(3)}(0)z^3+\cdots}
\EMdisp{\frac1{z-2}\EMbr =\frac{-1}{2-z}\EMbr =\frac{-1}2\,\frac1{1-z/2}\EMbr =\frac{-1}2\left(1+\frac z2+\frac{z^2}4+\frac{z^3}8+\cdots\right),\qquad\EMbr  \left|\frac z2\right|<1\ \EMbr \Longrightarrow\ |z|<2}
\EMdisp{\frac1{4+z}\EMbr =\frac14\,\frac1{1+z/4}\EMbr =\frac14\left(1-\frac z4+\frac{z^2}{16}-\frac{z^3}{64}+\cdots\right),\qquad\EMbr  \left|-\frac z4\right|<1\ \EMbr \Longrightarrow\ |z|<4}
\EMdisp{\frac d{dz}\left(\frac1{4+z}\right)\EMbr =\frac{-1}{(4+z)^2}\EMbr =\frac14\left(-\frac14+\frac{2z}{16}-\frac{3z^2}{64}+\cdots\right)\;\EMbr \Longrightarrow\;\frac1{(4+z)^2}\EMbr =\frac14\left(\frac14-\frac{2z}{16}+\frac{3z^2}{64}-\cdots\right)}
\EMdisp{f(z)\EMbr =(4z^2+30z+68)\times\frac1{(z+4)^2}\times\frac1{z-2}}
\EMdisp{f(z)\EMbr =(4z^2+30z+68)\times\left(\frac{-1}2\right)\left(1+\frac z2+\frac{z^2}4+\frac{z^3}8+\cdots\right)\times\frac14\left(\frac14-\frac z8+\frac{3z^2}{64}-\cdots\right)}
\EMdisp{f(z)\EMbr =\frac{-1}{16}\left(34+15z+\frac{67}8z^2+\cdots\right),\qquad\EMbr  |z|<2}
(region of convergence of radius 2.)

\end{EMexample}

\EMfigure{m06-radius-2}{The poles \EMim{z=-4} and \EMim{z=2}; the circle of convergence about
the origin has radius \EMim{2}}

\begin{EMexample}{}

Find the Laurent series of \EMim{f(z)=\cos z^2} in a neighbourhood of \EMim{z_0=0}.

From the definition of the cosine function,
\EMdisp{\cos z\EMbr =1-\frac{z^2}{2!}+\frac{z^4}{4!}-\frac{z^6}{6!}+-\cdots}
\EMdisp{f(z)\EMbr =\cos z^2\EMbr =1-\frac{z^4}{2!}+\frac{z^8}{4!}-\frac{z^{12}}{6!}+-\cdots}
The region of convergence is the whole complex plane, with infinite radius of convergence.

\end{EMexample}

\begin{EMexample}{}

Find the Laurent series of \EMim{f(z)=\dfrac{\cos z}{1-z^2}} in a neighbourhood of \EMim{z_0=0}.

\EMdisp{\left\{\begin{aligned}
\frac1{1-z^2}&=1+z^2+z^4+z^6+\cdots&&|z^2|<1\ \ \ |z|<1\\
\cos z&=1-\frac{z^2}{2!}+\frac{z^4}{4!}-\frac{z^6}{6!}+-\cdots&&|z|<\infty
\end{aligned}\right.}
\EMdisp{\begin{aligned}
f(z)=\frac{\cos z}{1-z^2}&=\big(1+z^2+z^4+z^6+\cdots\big)\left(1-\frac{z^2}{2!}+\frac{z^4}{4!}-\frac{z^6}{6!}+-\cdots\right)\\
&=1+\left(1-\frac1{2!}\right)z^2+\left(1+\frac1{4!}-\frac1{2!}\right)z^4+\cdots,\qquad |z|<1
\end{aligned}}

\end{EMexample}

\begin{EMexample}{}

Find the Laurent series of \EMim{f(z)=\dfrac1{z^4}} in a neighbourhood of \EMim{z_0=1}.

\EMdisp{g(z)\EMbr =\frac1z\EMbr =\frac1{z-1+1}\EMbr =\frac1{1+(z-1)}\EMbr =1-(z-1)+(z-1)^2-(z-1)^3+-\cdots,\qquad\EMbr  |z-1|<1}
\EMdisp{\frac d{dz}g(z)\EMbr =\frac{-1}{z^2}\EMbr =-1+2(z-1)-3(z-1)^2+4(z-1)^3-+\cdots,\qquad\EMbr  |z-1|<1}
\EMdisp{\frac{d^2}{dz^2}g(z)\EMbr =\frac2{z^3}\EMbr =2-6(z-1)+12(z-1)^2-20(z-1)^3+-\cdots,\qquad\EMbr  |z-1|<1}
\EMdisp{\frac{d^3}{dz^3}g(z)\EMbr =\frac{-6}{z^4}\EMbr =-6+24(z-1)-60(z-1)^2-+\cdots,\qquad\EMbr  |z-1|<1}
\EMdisp{f(z)\EMbr =\frac1{z^4}\EMbr =1-4(z-1)+10(z-1)^2-+\cdots}

\end{EMexample}

\EMfigure{m06-radius-3}{The region of convergence of the series about \EMim{z_0=1} is a disc of
radius \EMim{1} (up to the pole \EMim{z=0})}

\begin{EMexercise}{}

Find the Laurent expansions of the following functions in a neighbourhood of \EMim{z_0=0}:

\begin{enumerate}
\def\labelenumi{\arabic{enumi}.}
\tightlist
\item
  \EMim{f_1(z)=\dfrac1{\sqrt{1-z^2}}}
\item
  \EMim{f_2(z)=\sin^{-1}z}
\item
  \EMim{f_3(z)=\dfrac{e^z}{z^2(z^2+1)}}
\item
  Expand \EMim{\sinh z} in powers of \EMim{(z-\pi i)} and prove that \EMim{\displaystyle\lim_{z\to\pi i}\frac{\sinh z}{z-\pi i}=-1}
\item
  Prove that \EMim{\displaystyle z\cosh z^2=z+\sum_{n=1}^{\infty}\frac1{(2n)!}\,z^{4n+1}}
\end{enumerate}

\end{EMexercise}

\EMhold

\EMsection{Examples of Singular Points and Residues}{Examples of Singular Points and Residues}

\begin{EMexample}{}

Find the Laurent series of \EMim{f(z)=\dfrac{z}{(z+1)(z+4)^3}} in the neighbourhoods of its non-analytic points, and find the residue at those points.

The function is not analytic at \EMim{z=-1} and \EMim{z=-4}.

\textbf{(a)} In a neighbourhood of \EMim{z=-1}, a circle centred at this point with radius less than \EMim{3} (\EMim{|z+1|<3}) does not contain the isolated singular point \EMim{z=-4}.
\EMdisp{\frac{z}{(z+4)^3}\EMbr =\frac{z+4-4}{(z+4)^3}\EMbr =\frac{z+4}{(z+4)^3}+\frac{-4}{(z+4)^3}\EMbr =\frac1{(z+4)^2}-\frac4{(z+4)^3}}
\EMdisp{\frac1{z+4}\EMbr =\frac1{(z+1)+3}\EMbr =\frac13\,\frac1{1+\frac{z+1}3}\EMbr =\frac13\left[1-\frac{z+1}3+\left(\frac{z+1}3\right)^2-\left(\frac{z+1}3\right)^3+-\cdots\right]}
\EMdisp{\frac{-1}{(z+4)^2}\EMbr =\frac13\left[-\frac13+\frac2{3^2}(z+1)-\frac3{3^3}(z+1)^2+-\cdots\right]\;\EMbr \Longrightarrow\;\frac1{(z+4)^2}\EMbr =\frac19-\frac2{27}(z+1)+\frac1{27}(z+1)^2-+\cdots}
\EMdisp{\frac1{(z+4)^3}\EMbr =\frac16\left[\frac2{3^2}-\frac{2\times3}{3^3}(z+1)+\frac{3\times4}{3^4}(z+1)^2-\frac{4\times5}{3^5}(z+1)^3+-\cdots\right]\EMbr =\frac1{27}-\frac1{27}(z+1)+\frac2{81}(z+1)^2-+\cdots}
\EMdisp{\frac z{(z+4)^3}\EMbr =\frac1{(z+4)^2}-\frac4{(z+4)^3}\EMbr =-\frac1{27}+\frac2{27}(z+1)-\frac5{81}(z+1)^2+\cdots}
\EMdisp{f(z)\EMbr =\frac1{z+1}\,\frac z{(z+4)^3}\EMbr =-\frac1{27}\,\frac1{z+1}+\frac2{27}-\frac5{81}(z+1)+\cdots}
Type of singular point: a simple pole (of order \EMim{1}). Residue: \EMim{\operatorname{Res}\{f(z)\}\big|_{z=-1}=-\dfrac1{27}}.

\textbf{(b)} In a neighbourhood of \EMim{z=-4}, a circle centred at this point with radius less than \EMim{3} (\EMim{|z+4|<3}) does not contain the isolated singular point \EMim{z=-1}.
\EMdisp{\frac z{z+1}\EMbr =\frac{z+1-1}{z+1}\EMbr =\frac{z+1}{z+1}-\frac1{z+1}\EMbr =1-\frac1{z+1}}
\EMdisp{\frac1{z+1}\EMbr =\frac1{z+4-3}\EMbr =\frac1{-3+(z+4)}\EMbr =\frac{-1}3\,\frac1{1-\frac{z+4}3}\EMbr =\frac{-1}3\left[1+\frac{z+4}3+\frac{(z+4)^2}{3^2}+\cdots\right],\qquad\EMbr  |z+4|<3}
\EMdisp{\frac z{z+1}\EMbr =1+\frac13+\frac1{3^2}(z+4)+\frac1{3^3}(z+4)^2+\frac1{3^4}(z+4)^3+\cdots}
\EMdisp{f(z)\EMbr =\frac z{(z+1)(z+4)^3}\EMbr =\frac1{(z+4)^3}\,\frac z{z+1}\EMbr =\frac{4/3}{(z+4)^3}+\frac1{3^2}\frac1{(z+4)^2}+\frac1{3^3}\frac1{z+4}+\frac1{3^4}+\frac1{3^5}(z+4)+\frac1{3^6}(z+4)^2+\cdots}
Type of singular point: a pole of order \EMim{3}. Residue: \EMim{\operatorname{Res}\{f(z)\}\big|_{z=-4}=\dfrac1{27}}.

\end{EMexample}

\EMfigure{m06-poles-1}{The poles \EMim{z=-1} and \EMim{z=-4} of
\EMim{f(z)=\frac{z}{(z+1)(z+4)^3}}}

\begin{EMexample}{}

Find the Laurent series of \EMim{f(z)=e^{1/z}} in a neighbourhood of its non-analytic point, and find the residue at that point.

\EMdisp{e^z\EMbr =1+z+\frac{z^2}{2!}+\frac{z^3}{3!}+\cdots\;\EMbr \Longrightarrow\;e^{1/z}\EMbr =1+\frac1z+\frac1{2!\,z^2}+\frac1{3!\,z^3}+\cdots}
Type of singular point: an essential singularity (infinitely many negative powers). Residue: \EMim{\operatorname{Res}\{f(z)\}\big|_{z=0}=1}.

\end{EMexample}

\begin{EMexample}{}

Find the Laurent series of \EMim{f(z)=\cos\!\left(\dfrac z{z-1}\right)} in a neighbourhood of its non-analytic point, and find the residue at that point.

\EMdisp{\frac z{z-1}\EMbr =1+\frac1{z-1}\;\EMbr \Longrightarrow\;\cos\!\left(\frac z{z-1}\right)\EMbr =\cos\!\left(1+\frac1{z-1}\right)\EMbr =\cos1\cos\!\left(\frac1{z-1}\right)-\sin1\sin\!\left(\frac1{z-1}\right)}
\EMdisp{\cos z\EMbr =1-\frac{z^2}{2!}+\frac{z^4}{4!}-\frac{z^6}{6!}+-\cdots,\qquad\EMbr  \sin z\EMbr =z-\frac{z^3}{3!}+\frac{z^5}{5!}-+\cdots}
\EMdisp{\begin{aligned}
\cos\!\left(\frac z{z-1}\right)&=\cos1\left(1-\frac1{2!}\frac1{(z-1)^2}+\frac1{4!}\frac1{(z-1)^4}-+\cdots\right)\\
&\quad-\sin1\left(\frac1{z-1}-\frac1{3!}\frac1{(z-1)^3}+\frac1{5!}\frac1{(z-1)^5}-+\cdots\right)\\
&=\cos1-\sin1\,\frac1{z-1}-\cos1\,\frac1{2!}\frac1{(z-1)^2}+\sin1\,\frac1{3!}\frac1{(z-1)^3}+\cdots
\end{aligned}}
Type of singular point: an essential singularity. Residue: \EMim{\operatorname{Res}\{f(z)\}\big|_{z=1}=-\sin1}.

\end{EMexample}

\begin{EMtheorem}{Residue at a simple pole}

Prove that if \EMim{f(z)=q(z)/g(z)} and \EMim{g(z)} has a simple root at \EMim{z_0}, then
\EMdisp{\operatorname{Res}\{f(z)\}\Big|_{z=z_0}\EMbr =\frac{q(z_0)}{g'(z_0)}}

\end{EMtheorem}

\begin{EMproof}

Since \EMim{z_0} is a pole of order \EMim{1} we have (\EMim{p=1}):
\EMdisp{\begin{aligned}
\operatorname{Res}\{f(z)\}\Big|_{z=z_0}=A_1&=\frac1{(p-1)!}\left.\frac{d^{\,p-1}}{dz^{\,p-1}}\Big[(z-z_0)^pf(z)\Big]\right|_{z=z_0}=(z-z_0)\frac{q(z)}{g(z)}\Bigg|_{z=z_0}\\
&=(z-z_0)\frac{q(z)}{g(z)-g(z_0)}\Bigg|_{z=z_0}=\frac{q(z)}{\dfrac{g(z)-g(z_0)}{z-z_0}}\Bigg|_{z=z_0}=\frac{q(z_0)}{g'(z_0)}
\end{aligned}}

\end{EMproof}

\begin{EMexercise}{}

Find the Laurent series of the following functions about \EMim{z=0} and determine the type of this point.
\EMdisp{f_1(z)\EMbr =\frac z{e^z-1},\qquad\EMbr  f_2(z)\EMbr =\frac{e^{-z}}{z^3},\qquad\EMbr  f_3(z)\EMbr =\frac1{z^2(1-z^2)},\qquad\EMbr  f_4(z)\EMbr =\frac{\sin z^2}{z^3}}

\end{EMexercise}
